\chapter{Reinforcement learning from human feedback}
\label{ch:rlhf}

Large language models are trained first to predict text. To turn them into
helpful assistants, they are then fine-tuned with reinforcement learning,
with a reward that comes from human (or AI) judgements of their answers.
This last chapter shows that the tools of the course, PPO above all, apply
directly; and it derives direct preference optimisation, which removes the
RL loop altogether.

\section{A language model is a policy}

A language model $\pi_\theta(y \mid x)$ generates an answer $y = (y_1,
\dots, y_T)$ to a prompt $x$ token by token, with $\pi_\theta(y \mid x) =
\prod_t \pi_\theta(y_t \mid x, y_{<t})$. This is a policy in an MDP whose
state is the text so far, whose actions are the tokens of the vocabulary,
and whose transitions are deterministic (appending the token). The reward
is usually given only at the end, for the whole answer, so the problem can
also be seen as a contextual bandit (\cref{ch:bandits}) whose context is $x$
and whose ``arm'' is the whole answer $y$.

The standard pipeline (Ouyang et al., 2022, building on Christiano et al.,
2017) has three stages:
\begin{enumerate}[nosep]
  \item \textbf{Supervised fine-tuning:} behaviour cloning
    (\cref{ch:imitation}) on demonstrations of good answers, giving the
    reference policy $\pi_{\mathrm{ref}}$;
  \item \textbf{Reward modelling:} learning a reward $r_\phi(x, y)$ from
    comparisons between answers;
  \item \textbf{RL fine-tuning:} maximising $r_\phi$ with PPO, while staying
    close to $\pi_{\mathrm{ref}}$.
\end{enumerate}

\section{Reward models from preferences}

Humans compare answers more reliably than they score them. For a prompt
$x$, two answers are sampled and a labeller says which is better, $y_w$
(winner) or $y_l$ (loser). The Bradley--Terry model (1952) assumes that
\begin{equation}
  \label{eq:bt}
  \Prob(y_w \succ y_l \mid x) = \sigma\left( r(x, y_w) - r(x, y_l) \right),
  \qquad \sigma(u) = \frac{1}{1 + e^{-u}},
\end{equation}
for a latent reward $r$. The reward model, usually the language model with
a scalar head, is fitted by maximum likelihood:
\begin{equation}
  \label{eq:rm-loss}
  \mathcal{L}(\phi) = -\E_{(x, y_w, y_l)}\left[ \log \sigma\left( r_\phi(x,
  y_w) - r_\phi(x, y_l) \right) \right].
\end{equation}
The reward is only defined up to an additive function of $x$, which
cancels in \eqref{eq:bt}; this is harmless, since the policy is compared
on the same prompt.

\begin{intuition}
  This is the logic of Elo ratings in chess: each answer has a hidden
  score, and the probability that one is preferred to another depends only
  on the difference of their scores. Equal scores give $50\%$, a
  difference of $2$ gives $\sigma(2) \approx 88\%$. Fitting the model finds
  the scores that best explain which answers won their comparisons.
\end{intuition}

\section{RL fine-tuning with a KL penalty}

The policy is trained to maximise the learnt reward, regularised by a
Kullback--Leibler penalty towards the reference policy:
\begin{equation}
  \label{eq:rlhf-objective}
  \max_\pi\; \E_{x}\left[ \E_{y \sim \pi(\cdot \mid x)}[r(x, y)] - \beta\,
  \mathrm{KL}\left( \pi(\cdot \mid x) \,\|\, \pi_{\mathrm{ref}}(\cdot \mid
  x) \right) \right].
\end{equation}
The penalty matters: the reward model is accurate only near the answers
it was trained on, and an unconstrained policy finds answers where it is
wrong (\emph{reward hacking}, or overoptimisation: the true quality rises,
then falls as the learnt reward keeps rising). It also preserves the
fluency and the diversity of the reference model.

\begin{proposition}{Optimal policy of the KL-regularised objective}{rlhf-opt}
  For each prompt $x$, the objective \eqref{eq:rlhf-objective} is
  maximised by
  \begin{equation}
    \label{eq:rlhf-opt}
    \pi^\star(y \mid x) = \frac{1}{Z(x)}\, \pi_{\mathrm{ref}}(y \mid x) \exp
    \left( \frac{r(x, y)}{\beta} \right), \qquad Z(x) = \sum_y
    \pi_{\mathrm{ref}}(y \mid x)\, e^{r(x, y)/\beta}.
  \end{equation}
\end{proposition}

\begin{proof}
  For a fixed $x$ (omitted from the notation),
  \begin{align*}
    \E_\pi[r(y)] - \beta\, \mathrm{KL}(\pi \,\|\, \pi_{\mathrm{ref}})
    &= -\beta\, \E_{y \sim \pi}\left[ \log \frac{\pi(y)}{\pi_{\mathrm{ref}}
      (y) e^{r(y)/\beta}} \right] \\
    &= -\beta\, \mathrm{KL}(\pi \,\|\, \pi^\star) + \beta \log Z,
  \end{align*}
  maximal, since a KL divergence is non-negative and zero only between
  equal distributions, at $\pi = \pi^\star$.
\end{proof}

\begin{intuition}
  The optimal policy keeps the reference model's judgement of what is a
  plausible answer and tilts it towards rewarded answers: an answer whose
  reward is higher by $\beta$ becomes $e \approx 2.7$ times more likely. It
  cannot invent answers that the reference considers impossible ($\pi_{\mathrm
  {ref}} = 0$ stays $0$). The parameter $\beta$ sets how far the tilt may
  go, and so how much the learnt reward is trusted.
\end{intuition}

This is the same computation as soft policy improvement
(\cref{lem:soft-improvement}), with the reference policy as prior. The
normalising constant $Z(x)$ is a sum over all possible answers, so
\eqref{eq:rlhf-opt} cannot be used directly, and the objective is optimised
with PPO (\cref{ch:ppo}). In practice the KL penalty is spread over the
tokens: the reward of token $t$ is $-\beta \log(\pi_\theta(y_t \mid
\cdot) / \pi_{\mathrm{ref}}(y_t \mid \cdot))$, plus $r_\phi(x, y)$ at the
last token. A value head estimates the values, and GAE (\cref{def:gae})
computes the advantages; apart from the size of the networks, this is the
algorithm that trained the Pacman agent.

\paragraph{Without a critic.} The value network is as large as the policy,
which doubles the memory. Since the reward is given per answer, a baseline
can instead be computed by sampling several answers per prompt: GRPO (Shao
et al., 2024) samples $G$ answers, and uses as advantage of each answer its
reward standardised within the group, $A_i = (r_i - \mathrm{mean}(r)) /
\mathrm{std}(r)$, in the clipped PPO objective. It is the baseline of
\cref{ch:pg} estimated by Monte Carlo instead of a network. With rewards
that can be checked automatically (a mathematical answer, a program that
passes its tests), this RL with verifiable rewards trains models to
reason at length before answering.

\begin{intuition}
  GRPO grades each answer against the other attempts at the same question:
  an answer is reinforced if it did better than its siblings, discouraged
  if it did worse. The group plays the role of the critic, without a
  network: ``was this answer better than what I usually produce for this
  prompt?''
\end{intuition}

\section{Direct preference optimisation}

Rafailov et al.\ (2023) observed that \eqref{eq:rlhf-opt} can be inverted:
it expresses the reward in terms of the optimal policy,
\begin{equation}
  r(x, y) = \beta \log \frac{\pi^\star(y \mid x)}{\pi_{\mathrm{ref}}(y \mid
  x)} + \beta \log Z(x).
\end{equation}

\begin{theorem}{Direct preference optimisation}{dpo}
  Under the Bradley--Terry model \eqref{eq:bt}, with the reward
  parametrised by a policy as $r_\theta(x, y) = \beta \log (\pi_\theta(y
  \mid x) / \pi_{\mathrm{ref}}(y \mid x))$, the reward-model loss
  \eqref{eq:rm-loss} becomes
  \begin{equation}
    \label{eq:dpo}
    \mathcal{L}_{\mathrm{DPO}}(\theta) = -\E\left[ \log \sigma\left( \beta
    \log \frac{\pi_\theta(y_w \mid x)}{\pi_{\mathrm{ref}}(y_w \mid x)} -
    \beta \log \frac{\pi_\theta(y_l \mid x)}{\pi_{\mathrm{ref}}(y_l \mid x)}
    \right) \right],
  \end{equation}
  and $\pi_\theta$ is the optimal policy \eqref{eq:rlhf-opt} for the reward
  $r_\theta$. Minimising \eqref{eq:dpo} therefore fits a reward model and
  solves the KL-regularised RL problem for it at the same time.
\end{theorem}

\begin{proof}
  Adding $\beta \log Z(x)$ to $r_\theta$ does not change the
  Bradley--Terry probabilities, since it is the same for $y_w$ and $y_l$;
  hence the loss \eqref{eq:rm-loss} for $r_\theta$ is \eqref{eq:dpo}. For
  the reward $r_\theta$, the normalising constant of \eqref{eq:rlhf-opt} is
  $Z(x) = \sum_y \pi_{\mathrm{ref}}(y \mid x) \pi_\theta(y \mid x) /
  \pi_{\mathrm{ref}}(y \mid x) = 1$, and the optimal policy is
  $\pi_{\mathrm{ref}} \cdot \pi_\theta / \pi_{\mathrm{ref}} = \pi_\theta$.
  Conversely, any reward is equivalent, up to a function of $x$, to one of
  this form, namely the one of its optimal policy, so no generality is
  lost.
\end{proof}

\begin{intuition}
  In DPO the policy \emph{is} the reward model: the reward it implicitly
  assigns to an answer is how much more likely it makes that answer than
  the reference model does. Training raises this likelihood ratio for the
  preferred answer and lowers it for the rejected one. The gradient of
  \eqref{eq:dpo} is weighted by $\sigma(-\text{margin})$, so the pairs that
  the policy currently ranks wrongly count most, and those already ranked
  correctly with a large margin hardly count.
\end{intuition}

\begin{lstlisting}[caption={DPO loss (PyTorch); \code{logp} are summed token
  log-probabilities of the answers}]
def dpo_loss(logp_w, logp_l, ref_logp_w, ref_logp_l, beta=0.1):
    margin = beta * ((logp_w - ref_logp_w) - (logp_l - ref_logp_l))
    return -F.logsigmoid(margin).mean()
\end{lstlisting}

DPO is supervised learning on a fixed dataset of comparisons: no sampling
during training, no reward model, no value network. It is an
\emph{offline} method (\cref{ch:imitation}), and shares their limit: the
data are answers of another policy, whereas PPO with a reward model can
evaluate the new answers of the current policy. Online and iterative
variants sample fresh pairs from the current policy and label them.

\section{Feedback from AI}

Human labels are expensive and inconsistent. In RL from AI feedback
(RLAIF) the comparisons are made by a language model, guided by a written
set of principles (constitutional AI, Bai et al., 2022). Reward models can
also be trained to judge each step of a reasoning rather than the final
answer (process rewards). The difficulties are those this course met
throughout: the reward is a proxy, and an optimiser exploits its flaws;
the policy must explore to improve, but stay close to what is known to be
good; and the data distribution moves with the policy.

\section{Summary}

\begin{itemize}[nosep]
  \item A language model is a policy over tokens; RLHF fine-tunes it with a
    reward model learnt from pairwise comparisons (Bradley--Terry).
  \item The objective is the reward minus a KL penalty to the reference
    model; its optimum is the reference reweighted by $e^{r/\beta}$, and
    it is optimised with PPO and GAE, or with group baselines (GRPO).
  \item DPO inverts the optimal-policy formula and turns the problem into
    a classification loss on comparisons, with no RL loop.
\end{itemize}
